\begin{prob}
    Recall that the Bellman-Ford algorithm (with early stopping) will terminate early
    if, after updating every edge, no predecessors have changed. Suppose it is
    known that the greatest shortest path distance in a graph $G = (V, E)$ has
    $\Theta(\sqrt{V})$ edges. What is the worst case time complexity of
    Bellman-Ford when run on this graph? State your answer using $\Theta$
    notation.

    \begin{soln}
        The loop invariant for Bellman-Ford says that after updating all edges $\alpha$
        times, all nodes whose shortest path distance is $\leq \alpha$ has the correct
        shortest path distance. Since the longest shortest path is $\Theta(\sqrt{V})$,
        we'll find all shortest paths after $\Theta(\sqrt{V})$ (outer) iterations. Every
        iteration involves $\Theta(E)$ work because it updates every edge, so the total work
        of the loop is $\Theta(E \sqrt{V})$.

        We also spend $\Theta(V)$ time in the set up to initialize the estimated distance
        for each node in the graph. Therefore the total time is $\Theta(V + \sqrt V E)$.

        While $\Theta(V + \sqrt V E)$ gets full credit, we can technically
        simplify this even further. This is because we know that $E$ is at
        $\Omega(\sqrt{V})$, based on the fact that the longest shortest path
        has $\Theta(\sqrt{V})$ edges. Therefore, $E \sqrt{V}$ is $\Omega(V)$,
        so the $\sqrt{V} E$ term is at least as large as the $V$ term,
        asymptotically, and the time complexity can be simplified to
        $\Theta(\sqrt{V} E)$.
    \end{soln}
\end{prob}
